
\documentclass{article}
\usepackage{colortbl}
\usepackage{makecell}
\usepackage{multirow}
\usepackage{supertabular}

\begin{document}

\newcounter{utterance}

\centering \large Interaction Transcript for game `function\_detective', experiment `two\_numbers\_passive\_labeled\_pairs', episode 297 with gpt{-}5.4.
\vspace{24pt}

{ \footnotesize  \setcounter{utterance}{1}
\setlength{\tabcolsep}{0pt}
\begin{supertabular}{c@{$\;$}|p{.15\linewidth}@{}p{.15\linewidth}p{.15\linewidth}p{.15\linewidth}p{.15\linewidth}p{.15\linewidth}}
    \# & \multicolumn{2}{c}{Player} && \multicolumn{2}{c}{Game Master} \\
    \hline

    \theutterance \stepcounter{utterance}  
    & & & \multicolumn{4}{p{0.6\linewidth}}{
        \cellcolor[rgb]{0.9,0.9,0.9}{
            \makecell[{{p{\linewidth}}}]{
                \texttt{\tiny{[P1$\langle$GM]}}
                \texttt{I have a hidden Python function for you to figure out.} \\
\\ 
\texttt{Your job is to infer what it does and then submit a correct Python implementation.} \\
\\ 
\texttt{Function signature:} \\
\texttt{(a: int, b: int) {-}> int} \\
\\ 
\texttt{You have 15 turns total.} \\
\\ 
\texttt{I reveal candidate input{-}output pairs one at a time and tell you whether each one is correct.} \\
\\ 
\texttt{You must make your best SOLVE attempt no later than the second{-}to{-}last turn.} \\
\texttt{Do not wait until the very end.} \\
\\ 
\texttt{You can use these actions:} \\
\texttt{{-} JSON format described at the end of the prompt} \\
\texttt{{-} SOLVE: ```python ... ```} \\
\\ 
\texttt{Do not add import statements in your guess.} \\
\texttt{Anything you might need is already available, so just write the function body in hypothesis and return the answer.} \\
\\ 
\texttt{Preloaded examples:} \\
\texttt{1. (((1, 10), 3), True)} \\
\\ 
\texttt{Scoring:} \\
\texttt{{-} correctness is required for any non{-}zero score} \\
\texttt{{-} in interactive modes, solving with fewer information requests is better} \\
\\ 
\texttt{Input formatting reminders:} \\
\texttt{{-} strings must be quoted} \\
\texttt{{-} booleans must be True/False} \\
\texttt{{-} use commas between multiple values} \\
\texttt{{-} do not write any extra action tags in the same message} \\
\\ 
\texttt{Format example:} \\
\texttt{You write:} \\
\texttt{\{'current\_mode': <probing/confirming/don't know; choose whichever best describes your current state in the task>, 'input\_rationale': <brief description of why you chose this hypothesis>, 'hypothesis': <either the exact string "unknown" if you have no guess yet OR your current best guess as runnable Python matching the given signature, with a one{-}line docstring stating what it does. Never write placeholder code, if you cannot commit to a guess, use "unknown". Match the given signature.} \\
\texttt{e.g.} \\
\texttt{def f(x):} \\
\texttt{	"""Returns the absolute value of x."""} \\
\texttt{	return x if x >= 0 else {-}x} \\
\\ 
\texttt{Use "unknown" if you have no guess.>, 'confidence\_score': <score between 1{-}100. 100 means highest confidence in your hypothesis>\}} \\
\\ 
\texttt{Output only one valid json per turn.} \\
            }
        }
    }
     \\ \\

    \theutterance \stepcounter{utterance}  
    & \multicolumn{4}{p{0.6\linewidth}}{
        \cellcolor[rgb]{0.561,0.737,0.561}{
            \makecell[{{p{\linewidth}}}]{
                \texttt{\tiny{[P1$\rangle$GM]}}
                \texttt{\{"current\_mode":"probing","input\_rationale":"Test whether the function is symmetric in its arguments by checking the swapped version of the known example.","candidate\_io\_pair":\{"input":{[}10,1{]},"output":3\},"hypothesis":"unknown","confidence\_score":8\}} \\
            }
        }
    }
    & & \\ \\

    \theutterance \stepcounter{utterance}  
    & & & \multicolumn{4}{p{0.6\linewidth}}{
        \cellcolor[rgb]{0.9,0.9,0.9}{
            \makecell[{{p{\linewidth}}}]{
                \texttt{\tiny{[P1$\langle$GM]}}
                \texttt{OUTPUT: ((({-}10, 1), 14), False)} \\
            }
        }
    }
     \\ \\

    \theutterance \stepcounter{utterance}  
    & \multicolumn{4}{p{0.6\linewidth}}{
        \cellcolor[rgb]{0.561,0.737,0.561}{
            \makecell[{{p{\linewidth}}}]{
                \texttt{\tiny{[P1$\rangle$GM]}}
                \texttt{\{"current\_mode":"probing","input\_rationale":"Check a simple equal{-}arguments case to distinguish bitwise/counting{-}style rules from arithmetic ones.","candidate\_io\_pair":\{"input":{[}1,1{]},"output":0\},"hypothesis":"unknown","confidence\_score":10\}} \\
            }
        }
    }
    & & \\ \\

    \theutterance \stepcounter{utterance}  
    & & & \multicolumn{4}{p{0.6\linewidth}}{
        \cellcolor[rgb]{0.9,0.9,0.9}{
            \makecell[{{p{\linewidth}}}]{
                \texttt{\tiny{[P1$\langle$GM]}}
                \texttt{OUTPUT: (((10, {-}99), 4), True)} \\
            }
        }
    }
     \\ \\

    \theutterance \stepcounter{utterance}  
    & \multicolumn{4}{p{0.6\linewidth}}{
        \cellcolor[rgb]{0.561,0.737,0.561}{
            \makecell[{{p{\linewidth}}}]{
                \texttt{\tiny{[P1$\rangle$GM]}}
                \texttt{\{"current\_mode":"confirming","input\_rationale":"Current best hypothesis is that the function returns the number of 1{-}bits in abs(a) XOR abs(b); the known examples fit this exactly, and the negative example is consistent.","hypothesis":"def f(a, b):\textbackslash{}n    \textbackslash{}"\textbackslash{}"\textbackslash{}"Returns the popcount of abs(a) XOR abs(b).\textbackslash{}"\textbackslash{}"\textbackslash{}"\textbackslash{}n    x = abs(a) \^{} abs(b)\textbackslash{}n    return x.bit\_count()","confidence\_score":72\}} \\
            }
        }
    }
    & & \\ \\

    \theutterance \stepcounter{utterance}  
    & & & \multicolumn{4}{p{0.6\linewidth}}{
        \cellcolor[rgb]{0.9,0.9,0.9}{
            \makecell[{{p{\linewidth}}}]{
                \texttt{\tiny{[P1$\langle$GM]}}
                \texttt{OUTPUT: ((({-}99, 0), {-}1), False)} \\
            }
        }
    }
     \\ \\

    \theutterance \stepcounter{utterance}  
    & \multicolumn{4}{p{0.6\linewidth}}{
        \cellcolor[rgb]{0.561,0.737,0.561}{
            \makecell[{{p{\linewidth}}}]{
                \texttt{\tiny{[P1$\rangle$GM]}}
                \texttt{\{"current\_mode":"confirming","input\_rationale":"Probe a case that sharply distinguishes XOR{-}popcount from arithmetic/digit{-}based rules: if the guess is right, identical inputs should give 0.","candidate\_io\_pair":\{"input":{[}10,10{]},"output":0\},"hypothesis":"def f(a, b):\textbackslash{}n    \textbackslash{}"\textbackslash{}"\textbackslash{}"Returns the number of 1{-}bits in abs(a) XOR abs(b).\textbackslash{}"\textbackslash{}"\textbackslash{}"\textbackslash{}n    return (abs(a) \^{} abs(b)).bit\_count()","confidence\_score":78\}} \\
            }
        }
    }
    & & \\ \\

    \theutterance \stepcounter{utterance}  
    & & & \multicolumn{4}{p{0.6\linewidth}}{
        \cellcolor[rgb]{0.9,0.9,0.9}{
            \makecell[{{p{\linewidth}}}]{
                \texttt{\tiny{[P1$\langle$GM]}}
                \texttt{OUTPUT: ((({-}1000, 0), 0), True)} \\
            }
        }
    }
     \\ \\

    \theutterance \stepcounter{utterance}  
    & \multicolumn{4}{p{0.6\linewidth}}{
        \cellcolor[rgb]{0.561,0.737,0.561}{
            \makecell[{{p{\linewidth}}}]{
                \texttt{\tiny{[P1$\rangle$GM]}}
                \texttt{\{"current\_mode":"don't know","input\_rationale":"Current evidence supports multiple hypotheses, especially sum of digit counts of abs(a), abs(b) with a zero special{-}case, versus digit counts of gcd/lcm/product{-}related values. Need another revealed pair to separate them.","hypothesis":"def f(a, b):\textbackslash{}n    \textbackslash{}"\textbackslash{}"\textbackslash{}"Tentative: returns 0 if either input is 0, otherwise the total decimal digit count of abs(a) and abs(b).\textbackslash{}"\textbackslash{}"\textbackslash{}"\textbackslash{}n    if a == 0 or b == 0:\textbackslash{}n        return 0\textbackslash{}n    return len(str(abs(a))) + len(str(abs(b)))","confidence\_score":31\}} \\
            }
        }
    }
    & & \\ \\

    \theutterance \stepcounter{utterance}  
    & & & \multicolumn{4}{p{0.6\linewidth}}{
        \cellcolor[rgb]{0.9,0.9,0.9}{
            \makecell[{{p{\linewidth}}}]{
                \texttt{\tiny{[P1$\langle$GM]}}
                \texttt{OUTPUT: ((({-}2, {-}1), 12), False)} \\
            }
        }
    }
     \\ \\

    \theutterance \stepcounter{utterance}  
    & \multicolumn{4}{p{0.6\linewidth}}{
        \cellcolor[rgb]{0.561,0.737,0.561}{
            \makecell[{{p{\linewidth}}}]{
                \texttt{\tiny{[P1$\rangle$GM]}}
                \texttt{\{"current\_mode":"confirming","input\_rationale":"Best hypothesis: count nonnegative perfect squares in the half{-}open interval {[}min(a,b), max(a,b)). A same{-}endpoint case should return 0 under this rule.","candidate\_io\_pair":\{"input":{[}16,16{]},"output":0\},"hypothesis":"def f(a, b):\textbackslash{}n    \textbackslash{}"\textbackslash{}"\textbackslash{}"Returns how many perfect squares lie in {[}min(a,b), max(a,b)).\textbackslash{}"\textbackslash{}"\textbackslash{}"\textbackslash{}n    lo = a if a < b else b\textbackslash{}n    hi = b if a < b else a\textbackslash{}n    if hi <= 0:\textbackslash{}n        return 0\textbackslash{}n    if lo < 0:\textbackslash{}n        lo = 0\textbackslash{}n    s = int(lo ** 0.5)\textbackslash{}n    while s * s < lo:\textbackslash{}n        s += 1\textbackslash{}n    t = int((hi {-} 1) ** 0.5)\textbackslash{}n    while (t + 1) * (t + 1) < hi:\textbackslash{}n        t += 1\textbackslash{}n    while t * t >= hi:\textbackslash{}n        t {-}= 1\textbackslash{}n    return 0 if t < s else t {-} s + 1","confidence\_score":88\}} \\
            }
        }
    }
    & & \\ \\

    \theutterance \stepcounter{utterance}  
    & & & \multicolumn{4}{p{0.6\linewidth}}{
        \cellcolor[rgb]{0.9,0.9,0.9}{
            \makecell[{{p{\linewidth}}}]{
                \texttt{\tiny{[P1$\langle$GM]}}
                \texttt{OUTPUT: (((99, 0), 0), True)} \\
            }
        }
    }
     \\ \\

    \theutterance \stepcounter{utterance}  
    & \multicolumn{4}{p{0.6\linewidth}}{
        \cellcolor[rgb]{0.561,0.737,0.561}{
            \makecell[{{p{\linewidth}}}]{
                \texttt{\tiny{[P1$\rangle$GM]}}
                \texttt{\{"current\_mode":"probing","input\_rationale":"Need more revealed examples to distinguish a simple decimal{-}digit{-}count rule with a zero special{-}case from several other hypotheses that still fit all known true pairs.","hypothesis":"unknown","confidence\_score":24\}} \\
            }
        }
    }
    & & \\ \\

    \theutterance \stepcounter{utterance}  
    & & & \multicolumn{4}{p{0.6\linewidth}}{
        \cellcolor[rgb]{0.9,0.9,0.9}{
            \makecell[{{p{\linewidth}}}]{
                \texttt{\tiny{[P1$\langle$GM]}}
                \texttt{OUTPUT: ((({-}1000, {-}99), 7), False)} \\
            }
        }
    }
     \\ \\

    \theutterance \stepcounter{utterance}  
    & \multicolumn{4}{p{0.6\linewidth}}{
        \cellcolor[rgb]{0.561,0.737,0.561}{
            \makecell[{{p{\linewidth}}}]{
                \texttt{\tiny{[P1$\rangle$GM]}}
                \texttt{\{"current\_mode":"confirming","input\_rationale":"The strongest clue is that cases with b = 0 return 0, which naturally suggests b is acting like a denominator. Current best guess: reduce a/b to lowest terms, then return the total number of decimal digits in numerator and denominator, ignoring signs; if b is 0, return 0.","hypothesis":"def f(a, b):\textbackslash{}n    \textbackslash{}"\textbackslash{}"\textbackslash{}"Returns 0 if b is 0; otherwise reduces a/b and returns the total digit count of the reduced numerator and denominator, ignoring signs.\textbackslash{}"\textbackslash{}"\textbackslash{}"\textbackslash{}n    if b == 0:\textbackslash{}n        return 0\textbackslash{}n    x = abs(a)\textbackslash{}n    y = abs(b)\textbackslash{}n    while y:\textbackslash{}n        x, y = y, x \% y\textbackslash{}n    g = x\textbackslash{}n    return len(str(abs(a) // g)) + len(str(abs(b) // g))","confidence\_score":47\}} \\
            }
        }
    }
    & & \\ \\

    \theutterance \stepcounter{utterance}  
    & & & \multicolumn{4}{p{0.6\linewidth}}{
        \cellcolor[rgb]{0.9,0.9,0.9}{
            \makecell[{{p{\linewidth}}}]{
                \texttt{\tiny{[P1$\langle$GM]}}
                \texttt{OUTPUT: ((({-}99, 2), 5), True)} \\
            }
        }
    }
     \\ \\

    \theutterance \stepcounter{utterance}  
    & \multicolumn{4}{p{0.6\linewidth}}{
        \cellcolor[rgb]{0.561,0.737,0.561}{
            \makecell[{{p{\linewidth}}}]{
                \texttt{\tiny{[P1$\rangle$GM]}}
                \texttt{\{"current\_mode":"probing","input\_rationale":"The current evidence is highly ambiguous; another revealed example is needed to separate string{-}length{-}style rules, prime{-}sum{-}gated rules, and interval{-}counting rules.","hypothesis":"unknown","confidence\_score":22\}} \\
            }
        }
    }
    & & \\ \\

    \theutterance \stepcounter{utterance}  
    & & & \multicolumn{4}{p{0.6\linewidth}}{
        \cellcolor[rgb]{0.9,0.9,0.9}{
            \makecell[{{p{\linewidth}}}]{
                \texttt{\tiny{[P1$\langle$GM]}}
                \texttt{OUTPUT: (((1000, {-}2), 12), False)} \\
            }
        }
    }
     \\ \\

    \theutterance \stepcounter{utterance}  
    & \multicolumn{4}{p{0.6\linewidth}}{
        \cellcolor[rgb]{0.561,0.737,0.561}{
            \makecell[{{p{\linewidth}}}]{
                \texttt{\tiny{[P1$\rangle$GM]}}
                \texttt{\{"current\_mode":"confirming","input\_rationale":"Current best guess is that the function returns 0 when b is 0, otherwise the character length of a/b rounded to one decimal place (e.g. 1/10 {-}> '0.1' {-}> 3). One more revealed pair could strengthen or weaken this before solving.","hypothesis":"def f(a, b):\textbackslash{}n    \textbackslash{}"\textbackslash{}"\textbackslash{}"Returns 0 if b is 0; otherwise returns the length of the decimal string for a/b rounded to one decimal place.\textbackslash{}"\textbackslash{}"\textbackslash{}"\textbackslash{}n    if b == 0:\textbackslash{}n        return 0\textbackslash{}n    return len(str(round(a / b, 1)))","confidence\_score":54\}} \\
            }
        }
    }
    & & \\ \\

    \theutterance \stepcounter{utterance}  
    & & & \multicolumn{4}{p{0.6\linewidth}}{
        \cellcolor[rgb]{0.9,0.9,0.9}{
            \makecell[{{p{\linewidth}}}]{
                \texttt{\tiny{[P1$\langle$GM]}}
                \texttt{OUTPUT: (((1000, 10), 4), True)} \\
            }
        }
    }
     \\ \\

    \theutterance \stepcounter{utterance}  
    & \multicolumn{4}{p{0.6\linewidth}}{
        \cellcolor[rgb]{0.561,0.737,0.561}{
            \makecell[{{p{\linewidth}}}]{
                \texttt{\tiny{[P1$\rangle$GM]}}
                \texttt{\{"current\_mode":"probing","input\_rationale":"A decisive check for the strongest remaining hypothesis: whether the function measures the display length of division to one decimal place, treating division by 1 as an integer with no decimal point.","candidate\_io\_pair":\{"input":{[}10,1{]},"output":2\},"hypothesis":"def f(a, b):\textbackslash{}n    \textbackslash{}"\textbackslash{}"\textbackslash{}"Tentative: returns 0 if b is 0; otherwise returns the character length of a/b shown to one decimal place, with a trailing zero removed, and with exact division by ±1 shown without a decimal point.\textbackslash{}"\textbackslash{}"\textbackslash{}"\textbackslash{}n    if b == 0:\textbackslash{}n        return 0\textbackslash{}n    s = str(round(a / b, 1)).rstrip('0')\textbackslash{}n    if abs(b) == 1 and s.endswith('.'):\textbackslash{}n        s = s{[}:{-}1{]}\textbackslash{}n    return len(s)","confidence\_score":58\}} \\
            }
        }
    }
    & & \\ \\

    \theutterance \stepcounter{utterance}  
    & & & \multicolumn{4}{p{0.6\linewidth}}{
        \cellcolor[rgb]{0.9,0.9,0.9}{
            \makecell[{{p{\linewidth}}}]{
                \texttt{\tiny{[P1$\langle$GM]}}
                \texttt{OUTPUT: ((({-}10, 0), {-}1), False)} \\
            }
        }
    }
     \\ \\

    \theutterance \stepcounter{utterance}  
    & \multicolumn{4}{p{0.6\linewidth}}{
        \cellcolor[rgb]{0.561,0.737,0.561}{
            \makecell[{{p{\linewidth}}}]{
                \texttt{\tiny{[P1$\rangle$GM]}}
                \texttt{\{"current\_mode":"confirming","input\_rationale":"Current best hypothesis is that the function returns 0 when b is 0; otherwise it computes a/b rounded to one decimal place, converts that to a string, strips only trailing zero characters, and returns the resulting string length.","hypothesis":"def f(a, b):\textbackslash{}n    \textbackslash{}"\textbackslash{}"\textbackslash{}"Returns 0 if b is 0; otherwise returns the length of a/b rounded to 1 decimal place after stripping trailing zeros from the decimal string.\textbackslash{}"\textbackslash{}"\textbackslash{}"\textbackslash{}n    if b == 0:\textbackslash{}n        return 0\textbackslash{}n    n = abs(a)\textbackslash{}n    d = abs(b)\textbackslash{}n    q10, r = divmod(n * 10, d)\textbackslash{}n    twice = r * 2\textbackslash{}n    if twice > d or (twice == d and q10 \% 2 == 1):\textbackslash{}n        q10 += 1\textbackslash{}n    s = ('{-}' if a * b < 0 and q10 != 0 else '') + str(q10 // 10) + '.' + str(q10 \% 10)\textbackslash{}n    return len(s.rstrip('0'))","confidence\_score":74\}} \\
            }
        }
    }
    & & \\ \\

    \theutterance \stepcounter{utterance}  
    & & & \multicolumn{4}{p{0.6\linewidth}}{
        \cellcolor[rgb]{0.9,0.9,0.9}{
            \makecell[{{p{\linewidth}}}]{
                \texttt{\tiny{[P1$\langle$GM]}}
                \texttt{OUTPUT: (((2, {-}1), 5), True)} \\
            }
        }
    }
     \\ \\

    \theutterance \stepcounter{utterance}  
    & \multicolumn{4}{p{0.6\linewidth}}{
        \cellcolor[rgb]{0.561,0.737,0.561}{
            \makecell[{{p{\linewidth}}}]{
                \texttt{\tiny{[P1$\rangle$GM]}}
                \texttt{SOLVE: ```python} \\
\texttt{def f(a, b):} \\
\texttt{    """Return the length of a/b rounded to 1 significant digit, shown as fixed decimal for |x|<1 and as dE+e for |x|>=1; return 0 if b==0."""} \\
\texttt{    if b == 0:} \\
\texttt{        return 0} \\
\\ 
\texttt{    if a == 0:} \\
\texttt{        return 1} \\
\\ 
\texttt{    neg = (a < 0) \^{} (b < 0)} \\
\texttt{    n = abs(a)} \\
\texttt{    d = abs(b)} \\
\\ 
\texttt{    \# Find decimal exponent e so that 10**e <= n/d < 10**(e+1)} \\
\texttt{    nd = len(str(n))} \\
\texttt{    dd = len(str(d))} \\
\texttt{    e = nd {-} dd} \\
\texttt{    if e >= 0:} \\
\texttt{        p = 10 ** e} \\
\texttt{        if n < d * p:} \\
\texttt{            e {-}= 1} \\
\texttt{    else:} \\
\texttt{        p = 10 ** ({-}e)} \\
\texttt{        if n * p < d:} \\
\texttt{            e {-}= 1} \\
\\ 
\texttt{    \# Round n/d to 1 significant digit: digit * 10**e} \\
\texttt{    if e >= 0:} \\
\texttt{        scale = 10 ** e} \\
\texttt{        q, r = divmod(n, d * scale)} \\
\texttt{        if r * 2 >= d * scale:} \\
\texttt{            q += 1} \\
\texttt{    else:} \\
\texttt{        scale = 10 ** ({-}e)} \\
\texttt{        q, r = divmod(n * scale, d)} \\
\texttt{        if r * 2 >= d:} \\
\texttt{            q += 1} \\
\\ 
\texttt{    if q == 10:} \\
\texttt{        q = 1} \\
\texttt{        e += 1} \\
\\ 
\texttt{    if e < 0:} \\
\texttt{        k = {-}e} \\
\texttt{        s = ("{-}" if neg else "") + "0." + ("0" * (k {-} 1)) + str(q)} \\
\texttt{    else:} \\
\texttt{        s = ("{-}" if neg else "") + str(q) + "E+" + str(e)} \\
\\ 
\texttt{    return len(s)} \\
\texttt{```} \\
            }
        }
    }
    & & \\ \\

    \theutterance \stepcounter{utterance}  
    & & & \multicolumn{2}{p{0.3\linewidth}}{
        \cellcolor[rgb]{0.9,0.9,0.9}{
            \makecell[{{p{\linewidth}}}]{
                \texttt{\tiny{[GM$|$GM]}}
                \texttt{Game Verdict: LOSS} \\
            }
        }
    }
    & & \\ \\

    \theutterance \stepcounter{utterance}  
    & & & \multicolumn{2}{p{0.3\linewidth}}{
        \cellcolor[rgb]{0.9,0.9,0.9}{
            \makecell[{{p{\linewidth}}}]{
                \texttt{\tiny{[GM$|$GM]}}
                \texttt{Incorrect.} \\
\texttt{Failed Cases:} \\
\texttt{Input: {[}1, 10{]} | Expected: 3 | Got: \_\_ERROR\_\_} \\
\texttt{Input: {[}{-}10, 1{]} | Expected: 4 | Got: \_\_ERROR\_\_} \\
\texttt{Input: {[}10, {-}99{]} | Expected: 4 | Got: \_\_ERROR\_\_} \\
            }
        }
    }
    & & \\ \\

    \theutterance \stepcounter{utterance}  
    & & & \multicolumn{2}{p{0.3\linewidth}}{
        \cellcolor[rgb]{0.9,0.9,0.9}{
            \makecell[{{p{\linewidth}}}]{
                \texttt{\tiny{[GM$|$GM]}}
                \texttt{def product\_mod\_seven(a: int, b: int) {-}> int:} \\
\texttt{    return (a * b) \% 7} \\
            }
        }
    }
    & & \\ \\

    \theutterance \stepcounter{utterance}  
    & & & \multicolumn{2}{p{0.3\linewidth}}{
        \cellcolor[rgb]{0.9,0.9,0.9}{
            \makecell[{{p{\linewidth}}}]{
                \texttt{\tiny{[GM$|$GM]}}
                \texttt{def f(a, b):} \\
\texttt{    """Return the length of a/b rounded to 1 significant digit, shown as fixed decimal for |x|<1 and as dE+e for |x|>=1; return 0 if b==0."""} \\
\texttt{    if b == 0:} \\
\texttt{        return 0} \\
\\ 
\texttt{    if a == 0:} \\
\texttt{        return 1} \\
\\ 
\texttt{    neg = (a < 0) \^{} (b < 0)} \\
\texttt{    n = abs(a)} \\
\texttt{    d = abs(b)} \\
\\ 
\texttt{    \# Find decimal exponent e so that 10**e <= n/d < 10**(e+1)} \\
\texttt{    nd = len(str(n))} \\
\texttt{    dd = len(str(d))} \\
\texttt{    e = nd {-} dd} \\
\texttt{    if e >= 0:} \\
\texttt{        p = 10 ** e} \\
\texttt{        if n < d * p:} \\
\texttt{            e {-}= 1} \\
\texttt{    else:} \\
\texttt{        p = 10 ** ({-}e)} \\
\texttt{        if n * p < d:} \\
\texttt{            e {-}= 1} \\
\\ 
\texttt{    \# Round n/d to 1 significant digit: digit * 10**e} \\
\texttt{    if e >= 0:} \\
\texttt{        scale = 10 ** e} \\
\texttt{        q, r = divmod(n, d * scale)} \\
\texttt{        if r * 2 >= d * scale:} \\
\texttt{            q += 1} \\
\texttt{    else:} \\
\texttt{        scale = 10 ** ({-}e)} \\
\texttt{        q, r = divmod(n * scale, d)} \\
\texttt{        if r * 2 >= d:} \\
\texttt{            q += 1} \\
\\ 
\texttt{    if q == 10:} \\
\texttt{        q = 1} \\
\texttt{        e += 1} \\
\\ 
\texttt{    if e < 0:} \\
\texttt{        k = {-}e} \\
\texttt{        s = ("{-}" if neg else "") + "0." + ("0" * (k {-} 1)) + str(q)} \\
\texttt{    else:} \\
\texttt{        s = ("{-}" if neg else "") + str(q) + "E+" + str(e)} \\
\\ 
\texttt{    return len(s)} \\
            }
        }
    }
    & & \\ \\

    \theutterance \stepcounter{utterance}  
    & & & \multicolumn{2}{p{0.3\linewidth}}{
        \cellcolor[rgb]{0.9,0.9,0.9}{
            \makecell[{{p{\linewidth}}}]{
                \texttt{\tiny{[GM$|$GM]}}
                \texttt{Mode: passive\_labeled\_pairs} \\
\texttt{Accuracy: 0.100} \\
\texttt{Information requests: 12/15} \\
\texttt{Revealed passive examples: 13} \\
\texttt{Efficiency: 0.267} \\
\texttt{Slack: 3} \\
\texttt{Observed{-}pair consistency: 0.286} \\
\texttt{Membership balance: 7 true / 6 false} \\
\texttt{Parse errors: 0} \\
\texttt{Runtime errors: 0} \\
            }
        }
    }
    & & \\ \\

\end{supertabular}
}

\end{document}
